

Specifically, let \(\{(\mathcal{T}_k,\mathcal{V}_k)\}_{k=1}^{K}\) be
cross-validation splits of the epoch indices into training sets
\(\mathcal{T}_k\) and validation sets \(\mathcal{V}_k\). For a candidate
threshold \(\tau \in \mathcal{Q}_c\), the training epochs retained for channel
\(c\) in fold \(k\) are
\[
    \mathcal{G}_{k,c}(\tau)
    = \{ i \in \mathcal{T}_k : a_{i,c} \leq \tau \}.
\]
These retained epochs define a threshold-dependent training template
\[
    \bar{x}^{(\tau)}_{k,c}(t)
    = \frac{1}{|\mathcal{G}_{k,c}(\tau)|}
    \sum_{i \in \mathcal{G}_{k,c}(\tau)} x_{i,c,t},
\]
where candidates yielding an empty retained set are treated as inadmissible. The
corresponding validation reference is computed as the pointwise median across
the validation epochs,
\[
    \tilde{x}_{k,c}(t)
    = \operatorname{median}_{i \in \mathcal{V}_k} x_{i,c,t}.
\]
The use of a median validation reference makes the criterion less sensitive to
large validation-set artifacts. The validation loss for channel \(c\) is then
defined as
\[
    \mathcal{J}_{c}(\tau)
    = \frac{1}{K}\sum_{k=1}^{K}
    \left[
        \frac{1}{T}\sum_{t=1}^{T}
        \left\{\bar{x}^{(\tau)}_{k,c}(t)-\tilde{x}_{k,c}(t)\right\}^{2}
    \right]^{1/2}.
\]
The selected channel threshold is the candidate that minimizes the average  validation discrepancy,
\[
    \tau_c^\star
    = \arg\min_{\tau \in \mathcal{Q}_c} \mathcal{J}_{c}(\tau).
\]